\begin{afterword}
\summaryhead

The post argues that practical advice becomes more powerful when it is backed by experimental results, true causal accounts and valid mathematics. Its model is Seth Roberts, who lost weight on holiday while drinking unfamiliar-tasting drinks. Without the research on set points and flavor-calorie learning, the author says, Roberts would have written one blog post about a trick; with it, he built the Shangri-La Diet, which has helped many people though not the author.

The author's own advice to follow the evidence became clear in practice, the post says, only through probability theory. A theory also makes advice felt: a reader who knows about duration neglect and the peak-end rule sees that a large portion leaves almost the same memory as a small one, and saves dessert for last. The author wants advice on akrasia built the same way, and ranks the three kinds of backing by how hard they are to cite well.

\responsehead

The thesis is plausible, but what the post offers for it is illustration. It reports no comparison of advice given with and without a theory behind it, and its illustrations do not measure effects.

The lead example is the weakest on this count. The post says Roberts's knowledge of the science is ``the reason why as many people have benefited as they have.'' How many benefited was not measured. Fifteen days earlier the author had written that the diet's theory ``has never been analyzed by controlled experiment,'' and the evidence given there was reports on forums and blogs. The case shows that a theory led Roberts to a new procedure that many people tried and many reported success with.

The example of the dinner plate is the one that can be checked. The claim about portion size matches the one study of meals I found, published two years before the post: remembered pleasure did not depend on how much of a favorite dish people ate. The notes add what that study found about the aside on dessert.

Two points are information rather than criticism. The advice on akrasia the author asks for could already be written: Piers Steel's 2007 meta-analysis tied procrastination to hyperbolic discounting, and a LessWrong post built practical advice on it in 2011. And ego depletion, one of the topics the post lists, later failed a large replication, which fits the post's own caution that results at ``p < 0.05 may fail to replicate.''

In short: a plausible thesis supported by illustrations rather than evidence; the benefits in its lead example rest, by the author's own earlier account, on self-reports, and its point about portion size matches a study of the time.
\end{afterword}
